\section{Two free coordinates under an order-nine action}
\label{sec:e9-free-coordinates}

This section gives a conditional obstruction, not a forced-symmetry theorem.
An action by \emph{coordinate-preserving autotopies} of a group $E$ consists
of actions on the row, column and symbol sets $R,C,S$ such that
\[
 L(gr,gc)=gL(r,c)\qquad(g\in E).
\]
An action is \emph{free} if every point stabilizer is trivial. A view is
\emph{F} when all its two-line permutations have only even cycles.
The following elementary six-point observation supplies the obstruction.
Products of permutations act from right to left.

\begin{lemma}\label{lem:e9-six-point-product}
If $\beta,\gamma\in S_6$ both have cycle type $3+3$, then
$\gamma\beta^{-1}$ cannot be fixed-point-free with only even cycles.
\end{lemma}
\begin{proof}
Conjugate $\beta$ to $(0\,1\,2)(3\,4\,5)$. The unordered supports of
$\gamma$ are a partition into two triples. If they agree with those of
$\beta$, each restriction of $\gamma$ is that of $\beta$ or its inverse.
The product therefore has only cycles of lengths one or three.

Otherwise, exactly one support contains two points of $\{0,1,2\}$ and
one of $\{3,4,5\}$. The independent rotations of these two triples
centralize $\beta$ and act transitively on the nine possible mixed
partitions. We may consequently use supports $\{0,1,3\}$ and $\{2,4,5\}$.
Their four orientations give the following products:
\[
\begin{array}{c|c}
\gamma & \gamma\beta^{-1}\\ \hline
(0\,1\,3)(2\,4\,5)&(0\,4)(2\,3)(1)(5)\\
(0\,1\,3)(2\,5\,4)&(0\,5\,2\,3\,4)(1)\\
(0\,3\,1)(2\,4\,5)&(0\,4\,1\,3\,2)(5)\\
(0\,3\,1)(2\,5\,4)&(0\,5\,2)(1\,3\,4).
\end{array}
\]
Every case has a fixed point or an odd cycle. This covers all support
partitions and orientations, proving the assertion.
\end{proof}

\begin{theorem}\label{thm:e9-two-free-coordinates}
Let $L$ be a row-F Latin square of order $18$. Suppose
$E=C_3\times C_3$ acts on $L$ by coordinate-preserving autotopies,
freely on columns and symbols. Then its action on rows is also free.
\end{theorem}
\begin{proof}
Each free coordinate has two regular nine-point $E$-orbits. The possible
row-orbit sizes are $1,3,9$. We exclude the first two separately.

Suppose first that there is a row orbit $E/H$ of size three. The subgroup
$H$ of order three fixes all three of these rows because $E$ is abelian.
For each such row $r$, its bijection $f_r:C\to S$ is $H$-equivariant,
so it induces a bijection $\bar f_r:C/H\to S/H$ between six-point sets.
For two distinct selected rows, $p=f_r^{-1}f_t$ commutes with $H$.
Over a cycle of length $\ell$ of its induced permutation on $C/H$,
the return map $p^\ell$ is a translation of a regular three-point
$H$-orbit. Thus the physical cycles are either three cycles of length
$\ell$ or one cycle of length $3\ell$. In particular, an odd quotient
cycle would give an odd physical cycle. A quotient fixed point would
give a physical fixed point or a three-cycle, contrary respectively
to Latinness or row-F. Hence the three quotient rows form a row-F
$3\times6$ Latin rectangle. No column-F or symbol-F quotient is asserted.

The group $E/H$ rotates these three rows and has type $3+3$ on both
$C/H$ and $S/H$. For one generator, call the latter permutations
$\beta$ and $\gamma$. If $f$ is one quotient row, its successor is
$f_1=\gamma f\beta^{-1}$. The relative permutation is
\[
 f^{-1}f_1=(f^{-1}\gamma f)\beta^{-1}.
\]
Both factors have type $3+3$, contradicting
Lemma~\ref{lem:e9-six-point-product}. There is no size-three row orbit.

Finally, let $t$ be the number of globally $E$-fixed rows. Each such
row induces a bijection between the two regular column orbits and the
two regular symbol orbits. If two fixed rows induced the same bijection,
their relative map would preserve each regular column orbit and commute
with $E$. On a regular $E$-orbit it is a right translation. Since $E$
has exponent three, its cycles have length one or three, again impossible.
There are only two bijections of the orbit sets, so $t\leq2$. All other
row orbits have size divisible by three; therefore $t\equiv18\pmod3$,
giving $t=0$. All row orbits have size nine, as required.
\end{proof}

\begin{corollary}\label{cor:e9-three-view-freeness}
In an FFF Latin square of order $18$ with an actual
coordinate-preserving $C_3\times C_3$ action, freeness on any two
coordinate sets implies freeness on the third.
\end{corollary}
\begin{proof}
Permuting rows, columns and symbols transports the autotopy action and
FFF. Put the unspecified coordinate in the row position and apply
Theorem~\ref{thm:e9-two-free-coordinates}. This uses all three F views;
it does not assert arbitrary parastrophy invariance of row-F alone.
\end{proof}

\paragraph{Why both free coordinates are essential to this argument.}
Put $\alpha=(0\,3\,1)$, fixing $2,4,5$, and
$\beta=(0\,1\,2)(3\,4\,5)$. Then
$\alpha\beta^{-1}=(0\,2)(1\,3\,5\,4)$ has only even cycles.
More strongly, the following rectangle has row-F, with every relative
row map of type $4+2$:
\[
\begin{pmatrix}
0&1&2&3&4&5\\
2&4&0&1&5&3\\
5&2&1&0&3&4
\end{pmatrix}.
\]
Its rows $f_i$ satisfy $f_{i+1}(\alpha(c))=\beta(f_i(c))$, with row
indices modulo three. The column action has type $3+1+1+1$, outside
the lemma's hypotheses. This is a partial rectangle, not an FFF18 table.

\paragraph{Evidence and remaining scope.}
The proofs above are analytic. A separately implemented finite control
checks all $1600$ ordered pairs of type-$3+3$ factors using two product
conventions and two cycle calculations. It also checks $1600$ mixed
factor pairs, the displayed rectangle, literal quotient actions and
$18$ odd-fibre return maps. These checks are secondary to the proof,
not an enumeration of Latin squares. The elementary lemma and its
application carry no claim of literature priority. No assumption forces
an arbitrary FFF18 square to admit this symmetry: free actions, other
symmetry types and asymmetric squares are not decided. Unrestricted
order $18$ remains open.

The public control and its exact comparison target are in
\nolinkurl{manuscript/candidates/2026-10-02_e9_symmetry_review/}.
It can be rerun with Python 3.11+:
\begin{quote}\small
\begin{verbatim}
python3 -B tools/check_e9_symmetry.py \
  --output .audit/e9-symmetry-check
\end{verbatim}
\end{quote}
